% Auto-generated by extract-marking.py -- do not edit by hand unless you know what to do.
% Works for BOTH \documentclass[...,type=solutions,...] and
% [...,type=marking,...] compiles, chosen at compile time via
% \OlympMarkPick / \OlympMarkBeginTable (see olympiad-marking.sty).
% \eqref{n} falls back to a plain '(n)' whenever the referenced
% equation isn't actually typeset in the current document.

\OlympMarkBeginTable{|p{0.862\linewidth}|>{\centering\arraybackslash}p{0.098\linewidth}|}{|p{0.655\linewidth}|>{\centering\arraybackslash}p{0.085\linewidth}|>{\centering\arraybackslash}p{0.06\linewidth}|>{\centering\arraybackslash}p{0.06\linewidth}|}
\hline
\OlympMarkPick{\multicolumn{1}{|c|}{\textbf{Мазмұны}} & \textbf{Ұпайлар} \\ \hline}{\multicolumn{1}{|c|}{\textbf{Мазмұны}} & \textbf{Ұпайлар} & \textbf{MP} & \textbf{A} \\ \hline}
\OlympMarkPick{Формула \eqref{1}: $dP = \mu\cdot p dS\cdot \omega r,$ & 0.3 \\ \hline}{Формула \eqref{1}: $dP = \mu\cdot p dS\cdot \omega r,$ & 0.3 & &  \\ \hline}
\OlympMarkPick{\quad\textit{Балл $\omega r\Leftrightarrow v$ немесе $dS\Leftrightarrow2\pi rdr\Leftrightarrow rd\varphi dr$ немесе $pdS\Leftrightarrow dN$ жазбалары үшін де қойылады. $P=\mu N v$ немесе басқа \textbf{инфинитезималды емес} жазбалар үшін балл қойылмайды.} &  \\ \hline}{\quad\textit{Балл $\omega r\Leftrightarrow v$ немесе $dS\Leftrightarrow2\pi rdr\Leftrightarrow rd\varphi dr$ немесе $pdS\Leftrightarrow dN$ жазбалары үшін де қойылады. $P=\mu N v$ немесе басқа \textbf{инфинитезималды емес} жазбалар үшін балл қойылмайды.} &  & &  \\ \hline}
\OlympMarkPick{Формула \eqref{2}: $P = \frac23 \mu p\omega \pi R^3.$ & 0.5 \\ \hline}{Формула \eqref{2}: $P = \frac23 \mu p\omega \pi R^3.$ & 0.5 & &  \\ \hline}
\OlympMarkPick{\quad\textit{$2/3$ коэффициентінің орнына $1$ немесе $1/2$ жазылған. Егер \eqref{2} үшін \textbf{нөлдік емес} балл қойылса, онда \eqref{2}-дегі қателік \eqref{6}, \eqref{7} және \eqref{10}-да ескерілуі тиіс.} & \textit{-0.3} \\ \hline}{\quad\textit{$2/3$ коэффициентінің орнына $1$ немесе $1/2$ жазылған. Егер \eqref{2} үшін \textbf{нөлдік емес} балл қойылса, онда \eqref{2}-дегі қателік \eqref{6}, \eqref{7} және \eqref{10}-да ескерілуі тиіс.} & \textit{-0.3} & &  \\ \hline}
\OlympMarkPick{Формула \eqref{3}: $j = -\kappa \pi R^2\frac{dT}{dx},$ . Минус таңбасы міндетті емес, бірақ бұл \eqref{5}-тегі минус таңбасының болуына әсер етеді. & 0.3 \\ \hline}{Формула \eqref{3}: $j = -\kappa \pi R^2\frac{dT}{dx},$ . Минус таңбасы міндетті емес, бірақ бұл \eqref{5}-тегі минус таңбасының болуына әсер етеді. & 0.3 & &  \\ \hline}
\OlympMarkPick{Формула \eqref{4}: $dq = \alpha\cdot 2\pi R dx\cdot(T-T_0).$ & 0.3 \\ \hline}{Формула \eqref{4}: $dq = \alpha\cdot 2\pi R dx\cdot(T-T_0).$ & 0.3 & &  \\ \hline}
\OlympMarkPick{Формула \eqref{5}: $j(x+dx)-j(x)=dj=-dq.$ . Егер бұл \eqref{3} пен осы теңдеу арасындағы минус таңбаларының сәйкессіздігіне әкелсе, минус таңбасының болмауы \textbf{қабылданбайды}. & 0.3 \\ \hline}{Формула \eqref{5}: $j(x+dx)-j(x)=dj=-dq.$ . Егер бұл \eqref{3} пен осы теңдеу арасындағы минус таңбаларының сәйкессіздігіне әкелсе, минус таңбасының болмауы \textbf{қабылданбайды}. & 0.3 & &  \\ \hline}
\OlympMarkPick{Формула \eqref{6}: $\frac{d^2T}{dx^2} = \frac{2\alpha}{\kappa R}(T-T_0).$ & 0.2 \\ \hline}{Формула \eqref{6}: $\frac{d^2T}{dx^2} = \frac{2\alpha}{\kappa R}(T-T_0).$ & 0.2 & &  \\ \hline}
\OlympMarkPick{Формула \eqref{7}: $\lambda=\sqrt\frac{2\alpha}{\kappa R},$ & 0.2 \\ \hline}{Формула \eqref{7}: $\lambda=\sqrt\frac{2\alpha}{\kappa R},$ & 0.2 & &  \\ \hline}
\OlympMarkPick{Формула \eqref{8}: $T(x)=T_0+A\exp\ab(-\lambda|x|)$ & 0.4 \\ \hline}{Формула \eqref{8}: $T(x)=T_0+A\exp\ab(-\lambda|x|)$ & 0.4 & &  \\ \hline}
\OlympMarkPick{\quad\textit{Модуль міндетті емес. Егер $T_0$ орнына кез келген \textbf{нөлдік емес} тұрақты тұрса, қабылданады. Басқа жағдайларда балл қойылмайды.} &  \\ \hline}{\quad\textit{Модуль міндетті емес. Егер $T_0$ орнына кез келген \textbf{нөлдік емес} тұрақты тұрса, қабылданады. Басқа жағдайларда балл қойылмайды.} &  & &  \\ \hline}
\OlympMarkPick{Формула \eqref{9}: $j(0)=\frac P2.$ & 0.3 \\ \hline}{Формула \eqref{9}: $j(0)=\frac P2.$ & 0.3 & &  \\ \hline}
\OlympMarkPick{Формула \eqref{10}: $\omega = \frac3{\mu pR}\sqrt\frac{2\kappa\alpha }R\cdot(T^*-T_0)=\qty{51.0}{\radian\per\s}.$ & 0.5 \\ \hline}{Формула \eqref{10}: $\omega = \frac3{\mu pR}\sqrt\frac{2\kappa\alpha }R\cdot(T^*-T_0)=\qty{51.0}{\radian\per\s}.$ & 0.5 & &  \\ \hline}
\OlympMarkPick{Сандық жауабы \eqref{10} & 0.2 \\ \hline}{Сандық жауабы \eqref{10} & 0.2 & &  \\ \hline}
\OlympMarkPick{\quad\textit{Сандық есептеулерде қателіктің көшірілуі (error propagation) жоқ.} &  \\ \hline}{\quad\textit{Сандық есептеулерде қателіктің көшірілуі (error propagation) жоқ.} &  & &  \\ \hline}
\OlympMarkPick{\textbf{Барлығы} & \textbf{3.5} \\ \hline}{\textbf{Барлығы} & \textbf{3.5} &  &  \\ \hline}
\end{tabularx}
